% optimization/alignment/10_continuous_curvature_optimization.tex

\section{Formal Optimization in Curvature Space}

\subsection{Problem Setting}

Let \(L>0\) be fixed and let
\[
\kappa : [0,L] \to \R
\]
denote the curvature function of a transition.

We assume the boundary conditions
\[
\kappa(0)=0,
\qquad
\kappa(L)=\kappa_1,
\]
and optionally
\[
\kappa'(0)=0,
\qquad
\kappa'(L)=0.
\]

\subsection{Admissible Set}

\begin{definition}[Admissible transition set]
\[
\mathcal{A}
=
\left\{
\kappa \in H^1(0,L)
\mid
\kappa(0)=0,\;
\kappa(L)=\kappa_1
\right\}.
\]
\end{definition}

\subsection{Variational Problem}

\begin{definition}[Quadratic curvature-gradient functional]
\[
J[\kappa]
=
\int_0^L \bigl(\kappa'(s)\bigr)^2 \, \dd s .
\]
\end{definition}

\begin{definition}[Optimization problem]
Find $\kappa^\ast \in \mathcal{A}$ such that
\[
J[\kappa^\ast]
=
\min_{\kappa \in \mathcal{A}} J[\kappa].
\]
\end{definition}

\subsection{Euler--Lagrange Equation}

\begin{proposition}[Weak Euler--Lagrange equation]
\label{prop:curvature-el}
If \(\kappa^\ast\in\mathcal{A}\) minimizes \(J\) over \(\mathcal{A}\), then
\[
\kappa^{\ast\prime\prime}=0
\]
in the sense of distributions on \((0,L)\); in particular \(\kappa^\ast\)
is affine.
\end{proposition}

\begin{proof}
For every \(\eta\in H_0^1(0,L)\) and \(\epsilon\in\R\) the function
\(\kappa^\ast+\epsilon\eta\) lies in \(\mathcal{A}\). Minimality gives
\[
0=\frac{\dd}{\dd\epsilon}J[\kappa^\ast+\epsilon\eta]\Big|_{\epsilon=0}
=2\int_0^L\kappa^{\ast\prime}(s)\,\eta'(s)\,\dd s .
\]
Hence the distributional derivative of \(\kappa^{\ast\prime}\) vanishes,
so \(\kappa^{\ast\prime}\) is constant and \(\kappa^\ast\) is affine.
\end{proof}

\subsection{Interpretation}

\begin{proposition}[Unique minimizer]
\label{prop:curvature-el-minimizer}
The functional \(J\) has exactly one minimizer in \(\mathcal{A}\), namely
\[
\kappa^\ast(s)=\frac{\kappa_1}{L}\,s ,
\qquad
J[\kappa^\ast]=\frac{\kappa_1^2}{L}.
\]
\end{proposition}

\begin{proof}
By \cref{prop:curvature-el} a minimizer is affine; the boundary values
\(\kappa^\ast(0)=0\) and \(\kappa^\ast(L)=\kappa_1\) leave only
\(\kappa^\ast(s)=\kappa_1 s/L\). Existence and uniqueness also follow
directly from the Cauchy--Schwarz inequality: for every
\(\kappa\in\mathcal{A}\),
\[
\kappa_1^2
=\Bigl(\int_0^L\kappa'(s)\,\dd s\Bigr)^2
\le L\int_0^L\bigl(\kappa'(s)\bigr)^2\,\dd s
= L\,J[\kappa],
\]
with equality if and only if \(\kappa'\) is constant.
\end{proof}

Thus, the clothoid appears as the solution of a variational problem.
This connects classical railway transition design with an optimization
view in curvature space and provides a formal counterpart to the
historical role of the clothoid in railway engineering
\cite{Glover1900TransitionCurves,Higgins1922TransitionSpiral}.

\subsection{Higher-Order Model}

\begin{definition}
\[
J_2[\kappa]
=
\int_0^L \bigl(\kappa''(s)\bigr)^2 \,\dd s.
\]
\end{definition}

The functional \(J_2\) is not defined on the \(H^1\) set \(\mathcal{A}\);
it requires a second derivative and its own boundary conditions.

\begin{definition}[Clamped admissible set]
\[
\mathcal{A}_2
=
\left\{
\kappa \in H^2(0,L)
\mid
\kappa(0)=0,\;
\kappa(L)=\kappa_1,\;
\kappa'(0)=\kappa'(L)=0
\right\}.
\]
\end{definition}

\begin{proposition}[Fourth-order Euler--Lagrange equation]
\label{prop:curvature-el-fourth-order}
If \(\kappa^\ast\in\mathcal{A}_2\) minimizes \(J_2\) over \(\mathcal{A}_2\),
then \(\kappa^{\ast(4)}=0\) in the sense of distributions on \((0,L)\).
Consequently \(\kappa^\ast\) is the cubic polynomial
\[
\kappa^\ast(s)
=
\kappa_1\Bigl(3\bigl(\tfrac{s}{L}\bigr)^2-2\bigl(\tfrac{s}{L}\bigr)^3\Bigr),
\]
and it is the unique minimizer.
\end{proposition}

\begin{proof}
Admissible variations are \(\eta\in H_0^2(0,L)\), so that
\(\eta=\eta'=0\) at both ends. Stationarity gives
\(\int_0^L\kappa^{\ast\prime\prime}\eta''\,\dd s=0\) for all such
\(\eta\), i.e.\ \(\kappa^{\ast(4)}=0\) distributionally; hence
\(\kappa^\ast\) is a cubic polynomial. The four boundary conditions
determine its coefficients uniquely. The difference of two minimizers
lies in \(H_0^2(0,L)\) and has vanishing second derivative, hence is
zero; this gives uniqueness.
\end{proof}

Other boundary choices lead to natural boundary conditions instead. If
only \(\kappa(0)\) and \(\kappa(L)\) are prescribed, the natural
conditions are \(\kappa''(0)=\kappa''(L)=0\), and the minimizer reduces to
the affine law of \cref{prop:curvature-el-minimizer}. The fourth-order
equation itself is unchanged.

This explains the occurrence of cubic transition laws and links them to
modern polynomial transition design
\cite{ZboinskiWoznica2017Polynomial}.

\subsection{Generalized View}

General functionals of the form
\[
J[\kappa]
=
\int_0^L
F\bigl(\kappa,\kappa',\kappa'';v,c\bigr)\,\dd s
\]
allow incorporation of dynamic and engineering criteria.

\subsection{Connection to the Main Thesis}

\begin{summary}
Transition design can be formulated as an optimization problem in
curvature space.

Classical transition curves arise as solutions of such problems,
supporting the interpretation of transitions as curvature control
functions.
\end{summary}

\section{Coupled Optimization of Curvature and Cant}

\subsection{Motivation}

In railway applications, curvature and cant cannot be optimized
independently.
Both jointly determine the effective lateral acceleration and its rate
of change.

This motivates a coupled optimization problem in which both
\[
\kappa(s)
\quad \text{and} \quad
\phi(s)
\]
are treated as design variables.

Such coupling is fully consistent with railway alignment design
standards, where curvature, speed, and cant are not independent design
quantities \cite{EN13803_2017,Klein1998Trassierung}.

\subsection{Effective Acceleration Functional}

\begin{definition}[Effective lateral acceleration]
The effective lateral acceleration is
\[
a_{\mathrm{eff}}(s)
=
v^2 \kappa(s) - g \sin\phi(s).
\]
\end{definition}

A natural optimization objective is to reduce the magnitude of
unbalanced lateral acceleration.

\begin{definition}[Acceleration-based functional]
Define
\[
J_a[\kappa,\phi]
=
\int_0^L a_{\mathrm{eff}}(s)^2 \,\dd s
=
\int_0^L
\bigl(v^2 \kappa(s) - g \sin\phi(s)\bigr)^2
\,\dd s.
\]
\end{definition}

\subsection{Jerk-Based Functional}

\begin{definition}[Jerk]
The jerk is given by
\[
j(s)
=
v^3 \kappa'(s) - g v \cos\phi(s)\,\phi'(s).
\]
\end{definition}

\begin{definition}[Jerk functional]
Define
\[
J_j[\kappa,\phi]
=
\int_0^L j(s)^2 \,\dd s.
\]
That is,
\[
J_j[\kappa,\phi]
=
\int_0^L
\bigl(
v^3 \kappa'(s) - g v \cos\phi(s)\,\phi'(s)
\bigr)^2
\,\dd s.
\]
\end{definition}

This functional penalizes rapid variation of effective acceleration and
therefore provides a natural dynamic smoothness criterion. Such criteria
are closely related to engineering concerns about comfort and dynamic
compatibility \cite{EN13803_2017,BrustadDalmo2020Review}.

\subsection{Combined Objective}

\begin{definition}[Coupled objective functional]
A general coupled objective may be written as
\[
J[\kappa,\phi]
=
\alpha J_a[\kappa,\phi]
+
\beta J_j[\kappa,\phi]
+
\gamma J_\kappa[\kappa]
+
\delta J_\phi[\phi],
\]
where:
\[
J_\kappa[\kappa]
=
\int_0^L (\kappa'(s))^2 \,\dd s,
\qquad
J_\phi[\phi]
=
\int_0^L (\phi'(s))^2 \,\dd s.
\]
\end{definition}

The weights
\[
\alpha,\beta,\gamma,\delta \ge 0
\]
determine the relative importance of equilibrium, comfort, curvature
smoothness, and cant smoothness.

\subsection{Admissible Set}

\begin{definition}[Admissible railway state set]
Let
\[
\mathcal{B}
=
\left\{
(\kappa,\phi)
\;\middle|\;
\kappa \in H^1(0,L),\;
\phi \in H^1(0,L),
\;
\kappa(0)=0,\;
\kappa(L)=\kappa_1
\right\}.
\]
\end{definition}

Additional boundary conditions may be imposed, for example:
\[
\phi(0)=0,
\qquad
\phi(L)=\phi_1,
\]
or derivative constraints such as
\[
\kappa'(0)=\kappa'(L)=0,
\qquad
\phi'(0)=\phi'(L)=0.
\]

\subsection{Coupled Optimization Problem}

\begin{definition}[Coupled railway optimization problem]
Find
\[
(\kappa^\ast,\phi^\ast)\in\mathcal{B}
\]
such that
\[
J[\kappa^\ast,\phi^\ast]
=
\min_{(\kappa,\phi)\in\mathcal{B}} J[\kappa,\phi].
\]
\end{definition}

\subsection{Interpretation}

This problem extends classical transition design in two directions:
\begin{itemize}
\item from pure geometry to coupled geometry--cant design,
\item from curve selection to functional optimization.
\end{itemize}

In this framework, the transition curve is no longer chosen from a
catalogue, but emerges as part of an optimal coupled state.

\subsection{Special Cases}

\begin{corollary}[Substitution \(\phi\equiv 0\)]
\label{cor:coupled-substitution-phi-zero}
Substituting \(\phi\equiv 0\) removes cant as a variable. The coupled
functional reduces to the curvature-only functional
\[
J[\kappa,0]
=
\alpha v^4\int_0^L\kappa(s)^2\,\dd s
+
\beta v^6\int_0^L\bigl(\kappa'(s)\bigr)^2\,\dd s
+
\gamma J_\kappa[\kappa],
\]
which coincides with \(J_\kappa\) alone only for \(\alpha=\beta=0\).
\end{corollary}

\begin{corollary}[Balance identity]
\label{cor:balance-identity}
\(a_{\mathrm{eff}}(s)=0\) for all \(s\) if and only if
\[
v^2 \kappa(s)=g\sin\phi(s)
\qquad\text{for all } s .
\]
\end{corollary}

The identity is an algebraic consequence of the definition of
\(a_{\mathrm{eff}}\). Calling such a design \emph{perfectly balanced} is an
engineering interpretation that holds only under the constant-speed,
planar local approximation and the stated sign convention; it is not a
statement about vehicle dynamics.

These special cases show that classical equilibrium design appears as a
subcase of the coupled optimization framework.

\subsection{Relation to Optimal Control}

The coupled optimization problem may be interpreted as an optimal
control problem with controls
\[
\kappa(s), \quad \phi'(s),
\]
state variables
\[
x(s),y(s),\theta(s),\phi(s),
\]
and cost functional \(J[\kappa,\phi]\).

This provides a direct bridge from railway alignment design to modern
control-theoretic and variational methods.

\subsection{Connection to the Main Thesis}

\begin{summary}
The coupled optimization problem shows that railway alignment design is
not a purely geometric task.

Instead, it is the optimization of a coupled state system in which
curvature and cant jointly determine dynamic behaviour.

This supports the central thesis of the manuscript that transition
curves and cant evolution should be understood as control functions in a
curvature-based railway state model.
\end{summary}

\section{Summary}

\begin{summary}
Alignment optimization treats geometry as a variable rather than a fixed
input.

It combines parameterization, numerical reconstruction, and constrained
optimization methods to produce alignments that satisfy both geometric
and engineering requirements.
\end{summary}

